% --- Archon physics formalization source begin ---
% archon:physics
% archon:covers PhyXMiniProblems/problem_phyx_mini_0026.lean
% archon:source-report reports/phyx_mini/problem_phyx_mini_0026.source.json

\chapter{Physics problem phyx\_mini\_0026}
\label{ch:PhyXMiniProblems_problem_phyx_mini_0026}

\paragraph{Problem source.}
\#\# Physical scenario

A periscope as shown in figure is useful for viewing objects that cannot be seen directly. It can be used in submarines and when watching golf matches or parades from behind a crowd of people. Suppose the object is a distance \textbackslash{}( p\_1 \textbackslash{}) from the upper mirror and the centers of the two flat mirrors are separated by a distance \textbackslash{}( h \textbackslash{}).

\#\# Question

What is the distance of the final image from the lower mirror?

\#\# Answer choices

- A: A: \textbackslash{}( p\_1 - 2 h \textbackslash{})

- B: B: \textbackslash{}( p\_1 + 2 h \textbackslash{})

- C: C: \textbackslash{}( p\_1 - h \textbackslash{})

- D: D: \textbackslash{}( p\_1 + h \textbackslash{})

\#\# Figure caption (auxiliary; use the image as primary evidence)

The image illustrates a physics or optics diagram. It includes:

- A person standing to the left side, facing right, possibly holding a golf club.
- Two mirrors mounted vertically, facing each other. The mirrors are at a 45-degree angle inside a structure.
- A line labeled \textbackslash{}( p\_1 \textbackslash{}) extends horizontally from the person to the top mirror, indicating distance.
- A vertical line labeled \textbackslash{}( h \textbackslash{}) connects the top mirror to the bottom mirror, indicating height.
- An eye is depicted on the right side looking at the bottom mirror, indicating a line of sight or reflection observation.

\#\# Dataset metadata

Geometrical Optics; Spatial Relation Reasoning, Physical Model Grounding Reasoning

\paragraph{Recorded answer/context.}
D: D: \textbackslash{}( p\_1 + h \textbackslash{})

\paragraph{Figure/image path.}
/root/proposal\_for\_physic/science-mango/phyx\_mini\_run/phyx\_data/test\_image/26.png

\paragraph{Formalization target.}
create a compiling Lean file with sorry bodies at `PhyXMiniProblems/problem\_phyx\_mini\_0026.lean`. The Lean declarations must preserve the physical quantities, dimensions or dimensional roles, figure labels, governing-law hypotheses, and final relation expressed by this problem.
Use Mathlib/Physlib names found through LeanExplore where available. If a physics API is missing, introduce faithful local abstractions rather than scalar placeholder aliases.

\begin{theorem}[Physics formalization target]
\label{thm:physics:phyx_mini_0026:target}
\lean{PhyXMiniProblems.ProblemPhyXMini0026.problem_phyx_mini_0026}
\uses{def:physics:phyx-mini-0026:phyxminiproblems-problemphyxmini0026-metres, def:physics:phyx-mini-0026:phyxminiproblems-problemphyxmini0026-simetres, def:physics:phyx-mini-0026:phyxminiproblems-problemphyxmini0026-physicaldistance, def:physics:phyx-mini-0026:phyxminiproblems-problemphyxmini0026-periscopesetup, def:physics:phyx-mini-0026:phyxminiproblems-problemphyxmini0026-haspositivefigurelengths, def:physics:phyx-mini-0026:phyxminiproblems-problemphyxmini0026-matchesperiscopefigure, def:physics:phyx-mini-0026:phyxminiproblems-problemphyxmini0026-hasfortyfivedegreemirrorgeometry, def:physics:phyx-mini-0026:phyxminiproblems-problemphyxmini0026-satisfiesplanemirrorimageformation, def:physics:phyx-mini-0026:phyxminiproblems-problemphyxmini0026-finalimagematcheschoice}
The assigned autoformalize agent should translate this physics problem into Lean declarations in the covered file, with theorem and lemma proof bodies written as `by sorry`.
\end{theorem}
\begin{proof}
This is an autoformalization task, not a proof task. Produce statements that can later be proved by the physics prover without weakening the physical model.
\end{proof}

% --- Archon declaration topology begin ---
\section{Declaration topology}
\begin{definition}[Space]
  \label{def:physics:phyx-mini-0026:phyxminiproblems-problemphyxmini0026-space}
  \lean{PhyXMiniProblems.ProblemPhyXMini0026.Space}
  The Euclidean cross-section containing the object, mirrors, images, and eye
\end{definition}
\begin{proof}
  This is the declared model definition of Space.
\end{proof}
\begin{definition}[Length Quantity]
  \label{def:physics:phyx-mini-0026:phyxminiproblems-problemphyxmini0026-lengthquantity}
  \lean{PhyXMiniProblems.ProblemPhyXMini0026.LengthQuantity}
  A physical quantity carrying the dimension of length
\end{definition}
\begin{proof}
  This is the declared model definition of Length Quantity.
\end{proof}
\begin{definition}[metres]
  \label{def:physics:phyx-mini-0026:phyxminiproblems-problemphyxmini0026-metres}
  \lean{PhyXMiniProblems.ProblemPhyXMini0026.metres}
  \uses{def:physics:phyx-mini-0026:phyxminiproblems-problemphyxmini0026-lengthquantity}
  Turn an SI-metre scalar readout into a dimensionful physical length
\end{definition}
\begin{proof}
  This is the declared model definition of metres.
\end{proof}
\begin{definition}[si Metres]
  \label{def:physics:phyx-mini-0026:phyxminiproblems-problemphyxmini0026-simetres}
  \lean{PhyXMiniProblems.ProblemPhyXMini0026.siMetres}
  \uses{def:physics:phyx-mini-0026:phyxminiproblems-problemphyxmini0026-lengthquantity}
  Read a physical length as a real number of SI metres
\end{definition}
\begin{proof}
  This is the declared model definition of si Metres.
\end{proof}
\begin{definition}[physical Distance]
  \label{def:physics:phyx-mini-0026:phyxminiproblems-problemphyxmini0026-physicaldistance}
  \lean{PhyXMiniProblems.ProblemPhyXMini0026.physicalDistance}
  \uses{def:physics:phyx-mini-0026:phyxminiproblems-problemphyxmini0026-space, def:physics:phyx-mini-0026:phyxminiproblems-problemphyxmini0026-lengthquantity, def:physics:phyx-mini-0026:phyxminiproblems-problemphyxmini0026-metres}
  The dimensionful Euclidean distance between two points with SI-metre coordinates
\end{definition}
\begin{proof}
  This is the declared model definition of physical Distance.
\end{proof}
\begin{definition}[Plane Mirror]
  \label{def:physics:phyx-mini-0026:phyxminiproblems-problemphyxmini0026-planemirror}
  \lean{PhyXMiniProblems.ProblemPhyXMini0026.PlaneMirror}
  \uses{def:physics:phyx-mini-0026:phyxminiproblems-problemphyxmini0026-space}
  An ideal flat mirror in the two-dimensional periscope cross-section
\end{definition}
\begin{proof}
  This is the declared model definition of Plane Mirror.
\end{proof}
\begin{definition}[reflect Point]
  \label{def:physics:phyx-mini-0026:phyxminiproblems-problemphyxmini0026-planemirror-reflectpoint}
  \lean{PhyXMiniProblems.ProblemPhyXMini0026.PlaneMirror.reflectPoint}
  \uses{def:physics:phyx-mini-0026:phyxminiproblems-problemphyxmini0026-space, def:physics:phyx-mini-0026:phyxminiproblems-problemphyxmini0026-planemirror}
  The virtual image of a point in an ideal plane mirror
\end{definition}
\begin{proof}
  This is the declared model definition of reflect Point.
\end{proof}
\begin{definition}[Periscope Setup]
  \label{def:physics:phyx-mini-0026:phyxminiproblems-problemphyxmini0026-periscopesetup}
  \lean{PhyXMiniProblems.ProblemPhyXMini0026.PeriscopeSetup}
  \uses{def:physics:phyx-mini-0026:phyxminiproblems-problemphyxmini0026-space, def:physics:phyx-mini-0026:phyxminiproblems-problemphyxmini0026-lengthquantity, def:physics:phyx-mini-0026:phyxminiproblems-problemphyxmini0026-planemirror}
  All named physical objects and figure labels in the periscope diagram. horizontalAxis points toward the observer and verticalAxis points upward. The angles are real radian readouts. No field contains the requested final image distance
\end{definition}
\begin{proof}
  This is the declared model definition of Periscope Setup.
\end{proof}
\begin{definition}[Has Positive Figure Lengths]
  \label{def:physics:phyx-mini-0026:phyxminiproblems-problemphyxmini0026-haspositivefigurelengths}
  \lean{PhyXMiniProblems.ProblemPhyXMini0026.HasPositiveFigureLengths}
  \uses{def:physics:phyx-mini-0026:phyxminiproblems-problemphyxmini0026-simetres, def:physics:phyx-mini-0026:phyxminiproblems-problemphyxmini0026-periscopesetup}
  The two labeled distances are strictly positive physical lengths
\end{definition}
\begin{proof}
  This is the declared model definition of Has Positive Figure Lengths.
\end{proof}
\begin{definition}[Matches Periscope Figure]
  \label{def:physics:phyx-mini-0026:phyxminiproblems-problemphyxmini0026-matchesperiscopefigure}
  \lean{PhyXMiniProblems.ProblemPhyXMini0026.MatchesPeriscopeFigure}
  \uses{def:physics:phyx-mini-0026:phyxminiproblems-problemphyxmini0026-space, def:physics:phyx-mini-0026:phyxminiproblems-problemphyxmini0026-simetres, def:physics:phyx-mini-0026:phyxminiproblems-problemphyxmini0026-physicaldistance, def:physics:phyx-mini-0026:phyxminiproblems-problemphyxmini0026-periscopesetup}
  Geometric and scalar readouts supplied by the figure. In the chosen axes the object is p₁ to the left of the upper center, and the lower center is h below the upper center. The eye is on the positive, outgoing horizontal side
\end{definition}
\begin{proof}
  This is the declared model definition of Matches Periscope Figure.
\end{proof}
\begin{definition}[Has Forty Five Degree Mirror Geometry]
  \label{def:physics:phyx-mini-0026:phyxminiproblems-problemphyxmini0026-hasfortyfivedegreemirrorgeometry}
  \lean{PhyXMiniProblems.ProblemPhyXMini0026.HasFortyFiveDegreeMirrorGeometry}
  \uses{def:physics:phyx-mini-0026:phyxminiproblems-problemphyxmini0026-planemirror-reflectpoint, def:physics:phyx-mini-0026:phyxminiproblems-problemphyxmini0026-periscopesetup}
  The affine-reflection behavior of the two depicted 45-degree mirrors. The upper mirror sends any point a distance d left of its center to a virtual image the same distance above its center. The lower mirror sends any point a distance d above its center to a virtual image the same distance to its left. These universally quantified relations express mirror orientation and the equal-object/equal-image-distance law; neither mentions p₁ + h
\end{definition}
\begin{proof}
  This is the declared model definition of Has Forty Five Degree Mirror Geometry.
\end{proof}
\begin{definition}[Satisfies Plane Mirror Image Formation]
  \label{def:physics:phyx-mini-0026:phyxminiproblems-problemphyxmini0026-satisfiesplanemirrorimageformation}
  \lean{PhyXMiniProblems.ProblemPhyXMini0026.SatisfiesPlaneMirrorImageFormation}
  \uses{def:physics:phyx-mini-0026:phyxminiproblems-problemphyxmini0026-planemirror-reflectpoint, def:physics:phyx-mini-0026:phyxminiproblems-problemphyxmini0026-periscopesetup}
  Each successive virtual image is formed by reflection in the relevant mirror
\end{definition}
\begin{proof}
  This is the declared model definition of Satisfies Plane Mirror Image Formation.
\end{proof}
\begin{lemma}[first virtual image position]
  \label{lem:physics:phyx-mini-0026:phyxminiproblems-problemphyxmini0026-first-virtual-image-position}
  \lean{PhyXMiniProblems.ProblemPhyXMini0026.first_virtual_image_position}
  \uses{def:physics:phyx-mini-0026:phyxminiproblems-problemphyxmini0026-simetres, def:physics:phyx-mini-0026:phyxminiproblems-problemphyxmini0026-periscopesetup, def:physics:phyx-mini-0026:phyxminiproblems-problemphyxmini0026-haspositivefigurelengths, def:physics:phyx-mini-0026:phyxminiproblems-problemphyxmini0026-matchesperiscopefigure, def:physics:phyx-mini-0026:phyxminiproblems-problemphyxmini0026-hasfortyfivedegreemirrorgeometry, def:physics:phyx-mini-0026:phyxminiproblems-problemphyxmini0026-satisfiesplanemirrorimageformation}
  The first mirror places its virtual image p₁ above the upper mirror center
\end{lemma}
\begin{proof}
  The conclusion follows from the governing relations and figure readouts named in the statement of first virtual image position.
\end{proof}
\begin{lemma}[first virtual image position from lower]
  \label{lem:physics:phyx-mini-0026:phyxminiproblems-problemphyxmini0026-first-virtual-image-position-from-lower}
  \lean{PhyXMiniProblems.ProblemPhyXMini0026.first_virtual_image_position_from_lower}
  \uses{def:physics:phyx-mini-0026:phyxminiproblems-problemphyxmini0026-simetres, def:physics:phyx-mini-0026:phyxminiproblems-problemphyxmini0026-periscopesetup, def:physics:phyx-mini-0026:phyxminiproblems-problemphyxmini0026-haspositivefigurelengths, def:physics:phyx-mini-0026:phyxminiproblems-problemphyxmini0026-matchesperiscopefigure, def:physics:phyx-mini-0026:phyxminiproblems-problemphyxmini0026-hasfortyfivedegreemirrorgeometry, def:physics:phyx-mini-0026:phyxminiproblems-problemphyxmini0026-satisfiesplanemirrorimageformation}
  Because the lower mirror center is h below the upper center, the first virtual image is p₁ + h above the lower center
\end{lemma}
\begin{proof}
  The conclusion follows from the governing relations and figure readouts named in the statement of first virtual image position from lower.
\end{proof}
\begin{lemma}[final virtual image position]
  \label{lem:physics:phyx-mini-0026:phyxminiproblems-problemphyxmini0026-final-virtual-image-position}
  \lean{PhyXMiniProblems.ProblemPhyXMini0026.final_virtual_image_position}
  \uses{def:physics:phyx-mini-0026:phyxminiproblems-problemphyxmini0026-simetres, def:physics:phyx-mini-0026:phyxminiproblems-problemphyxmini0026-periscopesetup, def:physics:phyx-mini-0026:phyxminiproblems-problemphyxmini0026-haspositivefigurelengths, def:physics:phyx-mini-0026:phyxminiproblems-problemphyxmini0026-matchesperiscopefigure, def:physics:phyx-mini-0026:phyxminiproblems-problemphyxmini0026-hasfortyfivedegreemirrorgeometry, def:physics:phyx-mini-0026:phyxminiproblems-problemphyxmini0026-satisfiesplanemirrorimageformation}
  The lower mirror moves the final virtual image onto the horizontal output axis
\end{lemma}
\begin{proof}
  The conclusion follows from the governing relations and figure readouts named in the statement of final virtual image position.
\end{proof}
\begin{definition}[Answer Choice]
  \label{def:physics:phyx-mini-0026:phyxminiproblems-problemphyxmini0026-answerchoice}
  \lean{PhyXMiniProblems.ProblemPhyXMini0026.AnswerChoice}
  The four answer labels printed with the problem
\end{definition}
\begin{proof}
  This is the declared model definition of Answer Choice.
\end{proof}
\begin{definition}[answer Distance In Metres]
  \label{def:physics:phyx-mini-0026:phyxminiproblems-problemphyxmini0026-answerdistanceinmetres}
  \lean{PhyXMiniProblems.ProblemPhyXMini0026.answerDistanceInMetres}
  \uses{def:physics:phyx-mini-0026:phyxminiproblems-problemphyxmini0026-simetres, def:physics:phyx-mini-0026:phyxminiproblems-problemphyxmini0026-periscopesetup, def:physics:phyx-mini-0026:phyxminiproblems-problemphyxmini0026-answerchoice}
  The SI-metre expression printed beside each answer choice
\end{definition}
\begin{proof}
  This is the declared model definition of answer Distance In Metres.
\end{proof}
\begin{definition}[Final Image Matches Choice]
  \label{def:physics:phyx-mini-0026:phyxminiproblems-problemphyxmini0026-finalimagematcheschoice}
  \lean{PhyXMiniProblems.ProblemPhyXMini0026.FinalImageMatchesChoice}
  \uses{def:physics:phyx-mini-0026:phyxminiproblems-problemphyxmini0026-simetres, def:physics:phyx-mini-0026:phyxminiproblems-problemphyxmini0026-physicaldistance, def:physics:phyx-mini-0026:phyxminiproblems-problemphyxmini0026-periscopesetup, def:physics:phyx-mini-0026:phyxminiproblems-problemphyxmini0026-answerchoice, def:physics:phyx-mini-0026:phyxminiproblems-problemphyxmini0026-answerdistanceinmetres}
  A choice agrees exactly with the final-image distance along the central output ray
\end{definition}
\begin{proof}
  This is the declared model definition of Final Image Matches Choice.
\end{proof}
% --- Archon declaration topology end ---
% --- Archon physics formalization source end ---
